% Options for packages loaded elsewhere
\PassOptionsToPackage{unicode}{hyperref}
\PassOptionsToPackage{hyphens}{url}
%
\documentclass[
  a4paper,
]{book}

\usepackage{amsmath,amssymb}
\usepackage{iftex}
\ifPDFTeX
  \usepackage[T1]{fontenc}
  \usepackage[utf8]{inputenc}
  \usepackage{textcomp} % provide euro and other symbols
\else % if luatex or xetex
  \usepackage{unicode-math}
  \defaultfontfeatures{Scale=MatchLowercase}
  \defaultfontfeatures[\rmfamily]{Ligatures=TeX,Scale=1}
\fi
\usepackage{lmodern}
\ifPDFTeX\else  
    % xetex/luatex font selection
\fi
% Use upquote if available, for straight quotes in verbatim environments
\IfFileExists{upquote.sty}{\usepackage{upquote}}{}
\IfFileExists{microtype.sty}{% use microtype if available
  \usepackage[]{microtype}
  \UseMicrotypeSet[protrusion]{basicmath} % disable protrusion for tt fonts
}{}
\makeatletter
\@ifundefined{KOMAClassName}{% if non-KOMA class
  \IfFileExists{parskip.sty}{%
    \usepackage{parskip}
  }{% else
    \setlength{\parindent}{0pt}
    \setlength{\parskip}{6pt plus 2pt minus 1pt}}
}{% if KOMA class
  \KOMAoptions{parskip=half}}
\makeatother
\usepackage{xcolor}
\setlength{\emergencystretch}{3em} % prevent overfull lines
\setcounter{secnumdepth}{5}
% Make \paragraph and \subparagraph free-standing
\makeatletter
\ifx\paragraph\undefined\else
  \let\oldparagraph\paragraph
  \renewcommand{\paragraph}{
    \@ifstar
      \xxxParagraphStar
      \xxxParagraphNoStar
  }
  \newcommand{\xxxParagraphStar}[1]{\oldparagraph*{#1}\mbox{}}
  \newcommand{\xxxParagraphNoStar}[1]{\oldparagraph{#1}\mbox{}}
\fi
\ifx\subparagraph\undefined\else
  \let\oldsubparagraph\subparagraph
  \renewcommand{\subparagraph}{
    \@ifstar
      \xxxSubParagraphStar
      \xxxSubParagraphNoStar
  }
  \newcommand{\xxxSubParagraphStar}[1]{\oldsubparagraph*{#1}\mbox{}}
  \newcommand{\xxxSubParagraphNoStar}[1]{\oldsubparagraph{#1}\mbox{}}
\fi
\makeatother


\providecommand{\tightlist}{%
  \setlength{\itemsep}{0pt}\setlength{\parskip}{0pt}}\usepackage{longtable,booktabs,array}
\usepackage{multirow}
\usepackage{calc} % for calculating minipage widths
% Correct order of tables after \paragraph or \subparagraph
\usepackage{etoolbox}
\makeatletter
\patchcmd\longtable{\par}{\if@noskipsec\mbox{}\fi\par}{}{}
\makeatother
% Allow footnotes in longtable head/foot
\IfFileExists{footnotehyper.sty}{\usepackage{footnotehyper}}{\usepackage{footnote}}
\makesavenoteenv{longtable}
\usepackage{graphicx}
\makeatletter
\newsavebox\pandoc@box
\newcommand*\pandocbounded[1]{% scales image to fit in text height/width
  \sbox\pandoc@box{#1}%
  \Gscale@div\@tempa{\textheight}{\dimexpr\ht\pandoc@box+\dp\pandoc@box\relax}%
  \Gscale@div\@tempb{\linewidth}{\wd\pandoc@box}%
  \ifdim\@tempb\p@<\@tempa\p@\let\@tempa\@tempb\fi% select the smaller of both
  \ifdim\@tempa\p@<\p@\scalebox{\@tempa}{\usebox\pandoc@box}%
  \else\usebox{\pandoc@box}%
  \fi%
}
% Set default figure placement to htbp
\def\fps@figure{htbp}
\makeatother
% definitions for citeproc citations
\NewDocumentCommand\citeproctext{}{}
\NewDocumentCommand\citeproc{mm}{%
  \begingroup\def\citeproctext{#2}\cite{#1}\endgroup}
\makeatletter
 % allow citations to break across lines
 \let\@cite@ofmt\@firstofone
 % avoid brackets around text for \cite:
 \def\@biblabel#1{}
 \def\@cite#1#2{{#1\if@tempswa , #2\fi}}
\makeatother
\newlength{\cslhangindent}
\setlength{\cslhangindent}{1.5em}
\newlength{\csllabelwidth}
\setlength{\csllabelwidth}{3em}
\newenvironment{CSLReferences}[2] % #1 hanging-indent, #2 entry-spacing
 {\begin{list}{}{%
  \setlength{\itemindent}{0pt}
  \setlength{\leftmargin}{0pt}
  \setlength{\parsep}{0pt}
  % turn on hanging indent if param 1 is 1
  \ifodd #1
   \setlength{\leftmargin}{\cslhangindent}
   \setlength{\itemindent}{-1\cslhangindent}
  \fi
  % set entry spacing
  \setlength{\itemsep}{#2\baselineskip}}}
 {\end{list}}
\usepackage{calc}
\newcommand{\CSLBlock}[1]{\hfill\break\parbox[t]{\linewidth}{\strut\ignorespaces#1\strut}}
\newcommand{\CSLLeftMargin}[1]{\parbox[t]{\csllabelwidth}{\strut#1\strut}}
\newcommand{\CSLRightInline}[1]{\parbox[t]{\linewidth - \csllabelwidth}{\strut#1\strut}}
\newcommand{\CSLIndent}[1]{\hspace{\cslhangindent}#1}

\usepackage{titlesec}

\titleformat{\chapter}[display]
  {\normalfont\LARGE\bfseries} % Cambia \Huge por \LARGE para reducir el tamaño del título
  {\chaptertitlename\ \thechapter}{20pt}{\LARGE} % Ajusta el tamaño del título principal
\setlength{\headheight}{15pt} % Asegúrate de que este valor sea suficiente para el encabezado

\renewcommand{\chaptermark}[1]{%
  \markboth{\MakeUppercase{\truncate{50}{#1}}}{} % Limita la longitud del texto en el encabezado
}
\makeatletter
\@ifpackageloaded{bookmark}{}{\usepackage{bookmark}}
\makeatother
\makeatletter
\@ifpackageloaded{caption}{}{\usepackage{caption}}
\AtBeginDocument{%
\ifdefined\contentsname
  \renewcommand*\contentsname{Tabla de contenidos}
\else
  \newcommand\contentsname{Tabla de contenidos}
\fi
\ifdefined\listfigurename
  \renewcommand*\listfigurename{Listado de Figuras}
\else
  \newcommand\listfigurename{Listado de Figuras}
\fi
\ifdefined\listtablename
  \renewcommand*\listtablename{Listado de Tablas}
\else
  \newcommand\listtablename{Listado de Tablas}
\fi
\ifdefined\figurename
  \renewcommand*\figurename{Figura}
\else
  \newcommand\figurename{Figura}
\fi
\ifdefined\tablename
  \renewcommand*\tablename{Tabla}
\else
  \newcommand\tablename{Tabla}
\fi
}
\@ifpackageloaded{float}{}{\usepackage{float}}
\floatstyle{ruled}
\@ifundefined{c@chapter}{\newfloat{codelisting}{h}{lop}}{\newfloat{codelisting}{h}{lop}[chapter]}
\floatname{codelisting}{Listado}
\newcommand*\listoflistings{\listof{codelisting}{Listado de Listados}}
\makeatother
\makeatletter
\makeatother
\makeatletter
\@ifpackageloaded{caption}{}{\usepackage{caption}}
\@ifpackageloaded{subcaption}{}{\usepackage{subcaption}}
\makeatother

\ifLuaTeX
\usepackage[bidi=basic]{babel}
\else
\usepackage[bidi=default]{babel}
\fi
\babelprovide[main,import]{spanish}
% get rid of language-specific shorthands (see #6817):
\let\LanguageShortHands\languageshorthands
\def\languageshorthands#1{}
\usepackage{bookmark}

\IfFileExists{xurl.sty}{\usepackage{xurl}}{} % add URL line breaks if available
\urlstyle{same} % disable monospaced font for URLs
\hypersetup{
  pdftitle={Editar, publicar y financiar ciencia en América Latina},
  pdfauthor={Gimena del Rio Riande; Patricio Iván Pantaleo},
  pdflang={es},
  hidelinks,
  pdfcreator={LaTeX via pandoc}}


\title{Editar, publicar y financiar ciencia en América Latina}
\usepackage{etoolbox}
\makeatletter
\providecommand{\subtitle}[1]{% add subtitle to \maketitle
  \apptocmd{\@title}{\par {\large #1 \par}}{}{}
}
\makeatother
\subtitle{Perspectivas, experiencias y distopías en las dinámicas de la
comunicación científica en la región}
\author{Gimena del Rio Riande \and Patricio Iván Pantaleo}
\date{2025-01-16}

\begin{document}
\frontmatter
\maketitle

\renewcommand*\contentsname{Tabla de contenidos}
{
\setcounter{tocdepth}{2}
\tableofcontents
}

\mainmatter
\bookmarksetup{startatroot}

\chapter*{Preface}\label{preface}
\addcontentsline{toc}{chapter}{Preface}

\markboth{Preface}{Preface}

This is a Quarto book.

To learn more about Quarto books visit
\url{https://quarto.org/docs/books}.

Anna (2019)

\bookmarksetup{startatroot}

\chapter{O papel da Rede Ibero-Americana de Editores Científicos de
Educação (RIECE) no fortalecimento da comunicação científica em
periódicos da área de
Educação}\label{o-papel-da-rede-ibero-americana-de-editores-cientuxedficos-de-educauxe7uxe3o-riece-no-fortalecimento-da-comunicauxe7uxe3o-cientuxedfica-em-periuxf3dicos-da-uxe1rea-de-educauxe7uxe3o}

O capítulo trata acerca da fundação da Rede Ibero-americana de Editores
Científicos de Educação (RIECE) e do seu desenvolvimento de atividades
no âmbito da articulação internacional para o fortalecimento da
comunicação científica nos países latino-americanos, caribenhos e
africanos. O objetivo é apresentar à comunidade internacional como se
deu a articulação entre os editores dos países integrantes da rede para
a criação da RIECE destacando as experiências colaborativas exitosas
para majorar a internacionalização, indexação e aderência à ciência
aberta nos periódicos envolvidos. Metodologicamente, adota-se uma
pesquisa qualitativa, do tipo documental, que utiliza como fonte o
projeto da RIECE, o documento diagnóstico, as atas das reuniões e os
relatórios de atividades. Os resultados mostram que a RIECE se iniciou
com a articulação de 11 países - Brasil, Colômbia, México, Venezuela,
Costa Rica, Angola, Uruguai, Espanha, Portugal, Paraguai e Equador -- e
o apoio de importantes entidades que apoiam a editoria de revistas
científica, a exemplo do Redalyc, do Latindex, do EDUC@, da AURA e da
ABEC Brasil. Essa iniciativa foi aprovada pelo Conselho Nacional de
Desenvolvimento Científico e Tecnológico (CNPq), entidade ligada ao
Ministério da Ciência, Tecnologia e Inovações para incentivo à pesquisa
no Brasil, que concedeu apoio financeiro mediante verba de custeio e 15
bolsas internacionais, o que estimulou os participantes do projeto. O
trabalho inicial da RIECE amparou-se na pesquisa-ação, partindo de um
minucioso diagnóstico das revistas científicas, para, em seguida,
mediante questionário, grupo focal e relato textual elaborar um plano de
desenvolvimento colaborativo dos periódicos. Conclui-se que a RIECE vem
colaborando para qualificar científica e tecnologicamente os periódicos
de Educação a partir da articulação dos conhecimentos editoriais da
equipe compartilhados voluntariamente, consoante a ciência aberta, entre
os seus integrantes. Sua principal meta é ampliar o número de países da
rede e aumentar os editores de educação comprometidos com a proposta.

\hfill\break

\section{Introdução}\label{introduuxe7uxe3o}

O capítulo trata acerca da fundação da Rede Ibero-Americana de Editores
Científicos de Educação (RIECE) e do desenvolvimento de atividades no
âmbito da articulação internacional proporcionada pela RIECE para o
fortalecimento da comunicação científica nos países latino-americanos,
caribenhos e africanos.

A principal área coberta pela RIECE é a Educação, pois a articulação
centra-se em periódicos desse campo do conhecimento, com o mote de
qualificação contínua dos editores científicos considerando as
especificidades da área e a necessária articulação do sul global para o
fortalecimento da ciência nos países em desenvolvimento (Fialho et~al.,
2024). Ademais, a troca de saberes e conhecimentos editoriais entre os
editores tem sido fundamental para a formação profissional dos editores
científicos, já que, em quase a totalidade dos países latino-americanos,
caribenhos e africanos, não há uma habilitação específica em nível de
graduação ou pós-graduação para a formação de editores científicos,
inclusive, constata-se ausência de regulamentação para o exercício dessa
atividade laboral em diversos países (Morais et~al., 2022).

Majoritariamente, os editores de revistas científicas na Iberoamérica
são professores doutores, vinculados a instituições de ensino superior
públicas ou privadas, que assumem essa função como parte de suas
atividades laborais de pesquisa e gestão (Fontes, 2021). Destaca-se,
inclusive, que em diversos casos é realizado de maneira voluntária, sem
atribuição de carga horária ou benefício financeiro extra, em condições
precarizadas, pela ausência de uma política institucional adequada de
apoio e financiamento às necessidades da publicação científica (Werlang
et~al., 2021).

Considerando o exposto, o objetivo deste texto é apresentar à comunidade
internacional como se deu a articulação entre os editores de diversos
países para a criação da RIECE, destacando as experiências colaborativas
exitosas para majorar a internacionalização, indexação e aderência à
ciência aberta nos periódicos envolvidos, bem como a maneira pela qual
outros editores da área de Educação podem se juntar à Rede, somando
esforços para fortalecer as revistas científicas, no seu papel
fundamental de divulgar o conhecimento produzido por meio de pesquisas
qualificadas.

O Brasil, neste caso, partiu como idealizador do projeto em rede e,
inicialmente, coordenador central na articulação dos editores da área de
Educação, com o apoio financeiro do Conselho Nacional de Desenvolvimento
Científico e Tecnológico (CNPq), mediante aprovação de projeto
denominado ``Qualificação de periódicos científicos de Educação:
articular para internacionalizar'', submetido sob a coordenação da
pesquisadora brasileira Lia Machado Fiuza Fialho, na chamada do edital
publicizado pelo Ministério da Ciência, Tecnologia e Inovações (MCTI) e
CNPq, de número 14/2023. Outro apoio institucional importante foi o da
Fundação Cearense de Apoio ao Desenvolvimento Científico e Tecnológico
(Funcap), que também aprovou financiamento público para o projeto de
criação da RIECE mediante chamamento via edital público de número
06/2023. De tal modo, a iniciativa já começou com a aprovação e o
suporte financeiro de duas conceituadas instituições, respectivamente em
nível federal e estadual.

Consoante o edital MCTI/CNPq número 14/2023 e Funcap número 06/2023, o
projeto buscou atender ao objetivo de ``incentivar a participação de
pesquisadores brasileiros em projetos de cooperação com grupos e/ou
redes de pesquisa de países latino-americanos, caribenhos e africanos'';
bem como o objetivo de ``encorajar a participação em projetos
internacionais de pesquisadores brasileiros vinculados a instituições
sediadas nas regiões Norte, Nordeste e Centro-Oeste''. Dessa maneira,
inicialmente, para a constituição da Rede, foram priorizados editores de
países latino-americanos e das regiões brasileiras supramencionadas que
apresentam maior dificuldade de desenvolvimento científico, todavia sem
se olvidar de incluir países europeus com periódicos mais consolidados
que se sensibilizaram para ajudar na internacionalização dos demais.
Acrescenta-se ainda que, dentre os países mais desenvolvidos,
priorizaram-se os que falam os mesmos idiomas -- espanhol e português --
dos países latinos, com o objetivo de evitar barreiras linguísticas, já
que muitos editores não dominavam outros idiomas, a exemplo do francês e
do inglês.

Justamente por reconhecer a importância da RIECE para a articulação no
campo editorial ibero-americano e, em especial do sul global,
considerou-se relevante divulgar essa Rede para a comunidade
internacional, já que a propagação dessa iniciativa não apenas dar a ver
o trabalho previamente executado, mas possibilita a sua divulgação e o
seu fortalecimento mediante a união de outros editores da área de
Educação, o que permite à RIECE crescimento em número de adeptos e
fortalecimento para continuar atuando no fomento ao intercâmbio de
saberes e conhecimentos e à consolidação das revistas científicas.

\section{Metodologia}\label{metodologia}

Minayo (2008), ao debater sobre o conceito e a função da metodologia nas
pesquisas científicas, sinaliza para um enfoque plural, pois ``a
metodologia inclui as concepções teóricas de abordagem, o conjunto de
técnicas que possibilitam a apreensão da realidade e também o potencial
criativo do pesquisador'' (p.~22). Dessa maneira, é a partir do objetivo
da pesquisa que são determinadas as escolhas metodológicas do pesquisar
com vistas a eleger o melhor percurso científico para investigar a
questão proposta.

Para este estudo, adota-se uma pesquisa qualitativa, do tipo estudo de
caso único, que se utiliza de fonte documental, mais especificamente,
aquelas que dizem respeito ao projeto que resultou na criação RIECE,
aprovado pelo CNPq e pela Funcap, além do documento diagnóstico das
revistas inicialmente incluídas no projeto, das atas das reuniões e dos
relatórios de atividades. Essa escolha metodológica decorreu do fato de
que os documentos foram considerados as principais fontes para narrar a
história da fundação da RIECE.

A pesquisa é de abordagem qualitativa porque não está preocupada em
interpretar dados quantitativos e tecer generalizações (Lüdke \& André,
1986), ao contrário, busca compreender nuances subjetivas e condições
socioeconômicas e culturais que permitiram a articulação entre editores
da área de Educação para a criação da RIECE e o início do seu
desenvolvimento. Nessa direção, por tratar de uma Rede específica, a
investigação é do tipo estudo de caso único, já que centra ênfase na
análise de um caso particular e específico, o da RIECE (Yin, 2005).

As fontes documentais primárias foram eleitas para subsidiar a pesquisa
por serem consideradas importantes para compreender a história da RIECE.
De acordo Oliveira (2020), ``na pesquisa documental, o trabalho do(a)
pesquisador(a) requer uma análise mais cuidadosa, visto que os
documentos não passaram antes por nenhum tratamento científico''
(p.~70). A análise documental em tela, ao se utilizar de fontes
primárias, trabalhou com dados e informações que ainda não haviam sido
tratados cientificamente, possibilitando uma compreensão analítica
subjetiva e georreferenciada (Sá-Silva et~al., 2009).

Pimentel (2001) ressalta que no método de análise documental:

\begin{quote}
{[}\ldots{]} São descritos os instrumentos e meios de realização da
análise de conteúdo, apontando o percurso em que as decisões foram sendo
tomadas quanto às técnicas de manuseio de documentos: desde a
organização e classificação do material até a elaboração das categorias
de análise. (p.~179)
\end{quote}

Nessa direção, cabe informar que todos os documentos localizados que
diziam respeito à criação da RIECE foram considerados no estudo, sem
excluir nenhuma fonte. Projetos, atas e trabalhos apresentados e
publicados foram catalogados por data de confecção e analisados a partir
do tipo de fonte, ou seja, primeiro discutiu-se sobre o projeto inicial
e seu financiamento, que permitiu a criação da RIECE; em seguida, as
atas que registraram o desenvolvimento da Rede; e, por fim, as
atividades de divulgação e consolidação. Especificamente, as fontes
analisadas estão descritas adiante:

\begin{itemize}
\tightlist
\item
  Projeto submetido e aprovado pelo CNPq;
\item
  Projeto submetido e aprovado pela Funcap;
\item
  Termos de adesão assinados pelos integrantes do projeto;
\item
  Ata da reunião prévia para elaboração do projeto;
\item
  Ata da primeira reunião após aprovação do projeto;
\item
  Ata da segunda reunião do grupo com a criação da RIECE;
\item
  Material em Power Point de apresentação do projeto;
\item
  Anais de eventos com publicações sobre a RIECE;
\item
  Relatório diagnóstico das revistas envolvidas na RIECE;
\item
  Plataforma Carlos Chagas, onde se executa o projeto.
\end{itemize}

Todo o material consultado foi disponibilizado pela coordenadora do
projeto, que não apenas autorizou sua divulgação pública, por acreditar
que a ampla divulgação das pesquisas em todas as suas fases pode trazer
mais transparência ao processo investigativo, como também participou
ativamente da coleta de material e elaboração da redação do texto, com a
crença de que a possibilidade da divulgar a RIECE internacionalmente
pode ampliar e fortalecer essa iniciativa.

Cabe esclarecer que o projeto ``Qualificação de periódicos científicos
de Educação: articular para internacionalizar'', metodologicamente, para
alcançar os objetivos propostos, propôs-se a utilizar a pesquisa-ação
(Tripp, 2005), composta por cinco fases: 1) diagnóstica, com um
questionário misto para sondar os conhecimentos dos editores e
identificar as fortalezas e deficiências de cada periódico; 2)
planejamento da ação, com a elaboração do plano com atividades a serem
desenvolvidas para ampliar a internacionalização, indexação e aderência
à ciência aberta e minimizar as fragilidades identificadas; 3) tomada da
ação, com atividades práticas para a implementação de mudanças com
vistas a atender aos critérios dos principais indexadores, qualificar a
política editorial, alinhando-a à ciência aberta, e internacionalizar o
corpo científico de avaliadores e de autores; 4) avaliação, com grupo
focal, a partir da socialização das ações implementadas e sua discussão
circular; e 5) aprendizado em rede, com registro textual dos possíveis
benefícios e limites do trabalho cooperativo em rede internacional.
Inclusive, previa-se que as narrativas do grupo focal transcritas e dos
registros textuais seriam processadas no programa Interface de R pour
les Analyses Multidimensionnelles de Textes et de Questionnaires
(IRaMuTeQ) (Camargo \& Justo, 2018) com a análise de categoria para
posterior discussão. Todos os dados oriundos da pesquisa --
questionário, planejamento estratégico, relatório de mudanças realizadas
e registro textual -- também seriam depositados em acesso aberto no
LattesData. Por fim, o esperado é que os resultados possibilitem amparar
uma análise propositiva na elaboração de políticas públicas viáveis de
implementação pelas Instituições de Ensino Superior, Fundações de Amparo
à Pesquisa ou Ministérios de Educação dos países envolvidos.

Todavia, cabe salientar que o referido projeto ainda está em
desenvolvimento, de modo que o foco deste capítulo não são os resultados
do projeto ``Qualificação de periódicos científicos de Educação:
articular para internacionalizar'', e sim a criação da RIECE, um dos
produtos já alcançados pelo referido projeto, bem como as ações já
desenvolvidas pela Rede. Dessa maneira, após a análise documental, a
discussão será centrada na constituição da articulação inicial que
compôs a RIECE, ou seja, na sua fundação, bem como nas primeiras ações
da Rede em termos de diagnóstico das revistas participantes, da
participação em eventos e produções, e demais atividades desenvolvidas.

\section{Resultados e discussão}\label{resultados-e-discussuxe3o}

Os resultados, a partir da análise das fontes e termos de adesão
assinados pelos integrantes do projeto, mostraram que a RIECE se iniciou
com a articulação de 11 países -- Brasil, Colômbia, México, Venezuela,
Costa Rica, Angola, Uruguai, Espanha, Portugal, Paraguai e Equador -- e
o apoio de importantes entidades que dão suporte à editoria de revistas
científicas, a exemplo do Redalyc, do Latindex, do EDUC@, da Associação
Uruguaia de Revistas Científicas (AURA) e da ABEC Brasil.

Já a análise dos projetos interinstitucionais aprovados mostrou que essa
composição de 16 pesquisadores nacionais, dos quais seis são bolsistas
de produtividade em pesquisa do CNPq, com reconhecida experiência no
Brasil, e 16 pesquisadores internacionais, além de cinco representantes
de entidades do campo editorial (Redalyc, Latindex, EDUC@, AURA, ABEC),
foi impulsionada pelos critérios do edital de fomento, que priorizava a
participação do Brasil, de outros países do sul global, caribenhos e
africanos. De tal modo, o grupo envolveu 32 instituições, das quais 16
são nacionais e 16 são internacionais, com representação de todas as
cinco regiões do Brasil -- Norte, Nordeste, Sul, Sudeste e Centro-Oeste
- e o envolvimento de mais 11 países.

As instituições de vínculo dos pesquisadores integrantes deste projeto
estão apresentas no Tabla~\ref{tbl-1}, com a indicação da região
geográfica para as brasileiras e do país, no caso das internacionais.

\begin{longtable}[]{@{}ll@{}}
\caption{Instituições envolvidas no projeto}\label{tbl-1}\tabularnewline
\toprule\noalign{}
Instituições brasileiras & Região geográfica \\
\midrule\noalign{}
\endfirsthead
\toprule\noalign{}
Instituições brasileiras & Região geográfica \\
\midrule\noalign{}
\endhead
\bottomrule\noalign{}
\endlastfoot
Universidade Federal de Roraima & Norte \\
Universidade Federal de Rondônia & Norte \\
Universidade Federal do Oeste do Pará & Norte \\
Universidade Federal do Amazonas & Norte \\
Universidade Federal de Goiás & Centro-Oeste \\
Universidade Estadual de Goiás & Centro-Oeste \\
Universidade Federal de Catalão & Centro-Oeste \\
Universidade Federal do Vale do São Francisco & Nordeste \\
Universidade Federal da Paraíba & Nordeste \\
Universidade Federal do Ceará & Nordeste \\
Universidade Federal do Rio Grande do Norte & Nordeste \\
Centro Universitário UniChristus & Nordeste \\
Universidade Estadual do Ceará & Nordeste \\
Universidade Estadual de Ponta Grossa & Sul \\
Pontifícia Universidade Católica de São Paulo & Sudeste \\
Universidade Paulista Júlio de Mesquita Filho & Sudeste \\
\textbf{Instituições estrangeiras} & \textbf{País} \\
Universidad de Costa Rica & Costa Rica \\
Universidad de la República & Uruguai \\
Universidad Nacional Autónoma de México & México \\
Corporación Universitaria Minuto de Dios & Colômbia \\
Universidad Politécnica Salesiana & Equador \\
Universidad Central de Venezuela & Venezuela \\
Universidad Tecnológica Intercontinental & Paraguai \\
Universidad de Cádiz & Espanha \\
Universidad de Luanda Óscar Riba & Angola \\
Universidade de Coimbra & Portugal \\
Universidad de Valladolid & Espanha \\
Universidad Libre & Colômbia \\
Universidad Andina Simón Bolívar & Equador \\
Universidad Nacional Abierta & Venezuela \\
Universidad Católica del Uruguay & Uruguai \\
Universidad de Los Andes & Venezuela \\
\end{longtable}

A articulação dos editores decorreu do objetivo comum de aumentar a
qualificação dos periódicos em desenvolvimento, na área da Educação, a
partir da articulação internacional de editores científicos. Para isso,
buscava-se qualificar, científica e tecnologicamente, os periódicos
brasileiros e de outros países latino-americanos, caribenhos e africanos
a partir da articulação dos conhecimentos editoriais da equipe com foco
na internacionalização, na indexação e na aderência à ciência aberta.

Como no Tabla~\ref{tbl-1}, já foram apontadas as instituições envolvidas
no projeto, com suas respectivas regiões geográficas, no caso das
brasileiras, e dos países, no caso das estrangeiras, especificam-se
agora os membros da equipe com as revistas nas quais atuam como
editores, bem como seus respectivos grupos de pesquisas implicados no
projeto.

\begin{longtable}[]{@{}
  >{\raggedright\arraybackslash}p{(\linewidth - 4\tabcolsep) * \real{0.2625}}
  >{\raggedright\arraybackslash}p{(\linewidth - 4\tabcolsep) * \real{0.3625}}
  >{\raggedright\arraybackslash}p{(\linewidth - 4\tabcolsep) * \real{0.3750}}@{}}
\caption{Membros da equipe com a indicação das revistas científicas e
grupos de pesquisa a que estão vinculados}\label{tbl-2}\tabularnewline
\toprule\noalign{}
\begin{minipage}[b]{\linewidth}\raggedright
\textbf{Pesquisadores Internacionais}
\end{minipage} & \begin{minipage}[b]{\linewidth}\raggedright
\textbf{Revistas}
\end{minipage} & \begin{minipage}[b]{\linewidth}\raggedright
\textbf{Grupos de Pesquisa}
\end{minipage} \\
\midrule\noalign{}
\endfirsthead
\toprule\noalign{}
\begin{minipage}[b]{\linewidth}\raggedright
\textbf{Pesquisadores Internacionais}
\end{minipage} & \begin{minipage}[b]{\linewidth}\raggedright
\textbf{Revistas}
\end{minipage} & \begin{minipage}[b]{\linewidth}\raggedright
\textbf{Grupos de Pesquisa}
\end{minipage} \\
\midrule\noalign{}
\endhead
\bottomrule\noalign{}
\endlastfoot
Christian Andrés Torres Hurtado & Arista Crítica & Ágora
Latinoamericana \\
Benjamín Barón Velandia & Praxis Pedagógica & Laboratorio de Formación
Docente en Investigación \\
Cristina Maria Coimbra Vieira & Ex Aequo & Centro de Investigação em
Educação de Adultos e Intervenção Comunitária \\
Eva Álvarez Ramos & Ogigia & Literatura Española Contemporánea. Siglos
XX y XXI \\
Hugo Heredia Ponce & Hachetetepé. Revista Científica de Educación y
Comunicación & Investigación e Innovación Educativa en Didáctica de la
Lengua y la Literatura \\
Manuel Francisco Romero Oliva & Tavira & Investigación e Innovación
Educativa en Didáctica de la Lengua y la Literatura \\
Mariana Lacayo Campos & Educación & Publicaciones de la Revista
Educación \\
Tulio Ramírez & Areté & Problemas Educativos Venezolanos \\
Miguel Ángel Herrera Pavo & Revista Andina de Educación & Ciencia
Abierta Ecuador \\
Jaime Padilla Verdugo & Revista Científica Alteridad & Grupo de
Investigación en Ciencias de la Educación (GICCEE) \\
Eurico Wongo Gungula & Sapientiae & Grupo de Investigadores
Multidisciplinares sobre Ciencia Abierta \\
Rebeca Estefano de Salazar & Educ@ción en Contexto & Grupo de
Investigación de Educación \\
Héctor Magaña Vargas & Revista Mexicana de Orientación Educativa &
Orientación Educativa y Vocacional \\
Abelardo Juvenal Montiel Benitez & Tembikuaaty Rekávo & Grupo de
Investigación de la Universidad Tecnológica Intercontinental \\
María Alejandra Balbi & Páginas de Educación & Evaluación Formativa \\
Pedro José Rivas & Revista Venezolana en Educación & Grupo de
Investigación de Análisis Sociopolítico de Venezuela \\
\textbf{Pesquisadores nacionais} & \textbf{Revistas} & \textbf{Grupos de
pesquisa} \\
Lia Machado Fiuza Fialho & Educação \& Formação & Práticas Educativas,
Memórias e Oralidades \\
Charliton José dos Santos Machado & Ensino em Perspectivas & Grupo de
Estudos e Pesquisas História da Educação da Paraíba \\
Jefferson Mainardes & Práxis Educativa & Políticas Educacionais e
Práticas Educativas \\
Maria Lília Imbiriba Sousa Colares & Revista Exitu & História, Sociedade
e Educação no Brasil \\
Olivia Morais de Medeiros Neta & History of Education in Latin America &
Grupos de Pesquisa G-TRES \\
Cláudio Pinto Nunes & Revista Práxis Educacional & Grupo de Pesquisas
sobre Didática, Formação e Trabalho Docente \\
Alfrancio Ferreira Dias & Revista Tempos e Espaços em Educação & Grupo
de Estudos e Pesquisas Queer e Outras Epistemologias Feministas \\
Gilson Xavier Azevedo & Revista de Estudos em Educação & Educação e
Questões de Aprendizagem \\
José Gerardo Vasconcelos & Educação em Debates & História e Memória da
Educação \\
Marcelo Silva de Souza Ribeiro & Revista de Educação do Vale do São
Francisco & Educação e Desenvolvimento \\
Karla Angélica Silva do Nascimento & Educação \& Formação & Práticas
Educativas, Memórias e Oralidades \\
Fabiane Maia Garcia & Amazonida & Gênero, Trabalho e Educação \\
Rafael Fonseca de Castro & Práxis Pedagógica & Educação, Psicologia
Educacional e Processos Formativos \\
Karla Colares Vasconcelos & Revista Educação Pesquisa e Educação &
Formação de Professores Práticas Pedagógicas e Epistemologia do
Professor do Campo \\
Cristine Brandenburg & Práticas Educativas, Memórias e Oralidades &
Práticas Educativas, Memórias e Oralidades \\
Cláudia Tavares do Amaral & Poiésis Pedagógica & Centro de Investigação
e Estudos em Educação \\
Miriam Fábia Alves & Interações & Grupo de Estudos e Pesquisas em
Políticas Educacionais e Juventude \\
Francisca Genifer Andrade de Sousa & Práticas Educativas, Memórias e
Oralidades & Práticas Educativas, Memórias e Oralidades \\
\textbf{Pesquisador} & \textbf{Associação/Indexador} & \textbf{Grupos de
pesquisa} \\
Eduardo Aguado-López & Redalyc & Difusión y Divulgación del
Conocimiento \\
Ana María Cetto & Latindex & Grupo de Investigación Latindex \\
Nelson Antonio Simâo Gimenes & EDUC@ & Avaliação Educacional \\
Sylvia Laura Piovesan Suárez & AURA & Investigación CSIC Uruguay \\
Sigmar de Mello Rode & ABEC Brasil & Diagnóstico em Patologia Bucal \\
\end{longtable}

Importa salientar que os editores com periódicos mais qualificados se
dispuseram a ajudar na qualificação dos periódicos emergentes,
especialmente no tocante à indexação, à internacionalização e à
adequação à ciência aberta. Somando-se a essa articulação prévia que
mobiliza 11 países e 34 editores científicos da área de Educação, a
ideia do projeto foi previamente aprovada e apoiada por cinco
representantes de entidades muito importantes no campo editorial: duas
das mais consolidadas associações científicas de editores dos países
latino-americanos, a AURA e a ABEC Brasil; bem como pelos indexadores
Redalyc, Latindex e EDUC@, de modo que essas entidades assumiram o
compromisso de colaborar com o projeto seja no âmbito da sua divulgação,
da formação dos editores ou da assessoria para qualificação,
internacionalização e indexação dos periódicos.

O projeto já se iniciou robusto em colaborações concretas de
internacionalização, com a proposta coletiva entre os membros envolvidos
de iniciar uma rede permanente de colaboração envolvendo países
latino-americanos que se amplie e se perpetue mesmo após a finalização
do projeto. Nesse sentido, a rede colaborativa aqui descrita, a RIECE,
apresentou forte potencial de se converter em uma consolidada rede de
editores da área de Educação na Iberoamérica.

Previam-se alguns objetivos específicos: 1) fortalecer o desenvolvimento
tecnológico de periódicos científicos em Educação desde o fortalecimento
da participação de pesquisadores brasileiros e estrangeiros em projeto
de cooperação em rede internacional; 2) incentivar a participação de
pesquisadores nacionais e internacionais, editores de revistas
científicas, vinculados às instituições localizadas geograficamente nas
regiões menos favorecidas, em projetos internacionais com países
latino-americanos e caribenhos que enfrentam problemas semelhantes para
o aprimoramento das revistas científicas; 3) desenvolver projeto de
extensão, envolvendo jovens doutores, para a promoção de ações de
formação no campo editorial (para editores e autores) com vistas a
ampliar também a articulação com a comunidade, alcançando também o
público leigo; 4) promover maior internacionalização, divulgação
científica qualificada, aderência à ciência aberta, indexação e boas
práticas editoriais nas revistas científicas nacionais e internacionais,
em especial, nas envolvidas diretamente com o projeto; e 5) constituir
uma rede internacional de editores científicos parceiros com o objetivo
comum de aprimorar o trabalho editorial e a qualificação dos periódicos
em Educação do sul global.

Observou-se que, mesmo ainda em fase quase embrionária, com menos de um
ano do seu início formal, já se conseguia alcançar todos os objetivos
propostos, afinal foi possível articular editores de diversos países,
constituindo uma rede de cooperação internacional, a RIECE, que envolvia
em projeto internacional países latino-americanos e caribenhos. Ademais,
já se desenvolviam iniciativas de promoção de ações de formação no campo
editorial, inclusive fortalecendo a cultura da ciência aberta e a
aproximação com a comunidade na perspectiva da ciência cidadã, bem como
já havia se realizado um sério diagnóstico das revistas envolvidas na
Rede para majorar a internacionalização, aderência à ciência aberta,
boas práticas editoriais e indexação das revistas.

A ata da reunião prévia para a elaboração do projeto mostrou que a
iniciativa partiu de uma articulação de editores brasileiros em
Educação, os quais se reuniram para discutir o edital do CNPq, informar
da impossibilidade legal de estar concorrendo em mais de uma proposta,
assumir o compromisso com as atividades da pesquisa-ação previstas para
o desenvolvimento dos periódicos e elaborar coletivamente o orçamento,
especialmente no que dizia respeito às bolsas para missões de
internacionalização nos países parceiros, já que era necessário constar
no corpo do texto da proposta o nome dos brasileiros interessados em
trocar saberes e experiências com editores de outros periódicos de
Educação de países distintos.

As bolsas foram divididas equitativamente entre todos os membros
brasileiros que manifestaram interesse, de modo que restou uma bolsa
para Pós-Doutorado Junior no Exterior (PDE), 13 bolsas de
Desenvolvimento Tecnológico e Inovação no Exterior na categoria Sênior
(DES) e uma bolsa de Desenvolvimento Tecnológico e Inovação no Exterior
na categoria Junior (DEJ).

\begin{longtable}[]{@{}lll@{}}
\caption{Quantitativo de bolsas financiadas}\label{tbl-3}\tabularnewline
\toprule\noalign{}
Tipo de bolsa & Descrição & Quantidade \\
\midrule\noalign{}
\endfirsthead
\toprule\noalign{}
Tipo de bolsa & Descrição & Quantidade \\
\midrule\noalign{}
\endhead
\bottomrule\noalign{}
\endlastfoot
PDE & 6 meses & 1 \\
DEJ & 1 mês & 1 \\
DES & 1 mês & 13 \\
\end{longtable}

A maioria das bolsas eram de categoria Sênior em decorrência de a equipe
possuir apenas um jovem doutor com menos de cinco anos da conclusão de
seu curso de doutorado e todos os demais já serem pesquisadores
experientes. Esse perfil era esperado, tendo em vista que, para assumir
a função de editor de revista científica, geralmente são eleitos
pesquisadores com maior visibilidade, inserção e maturidade com a
produção e divulgação do conhecimento (Anna, 2019; Gomes, 2010).

A ata referente à primeira reunião após aprovação do projeto mostra que
o objetivo geral do encontro foi apresentar: os membros da equipe
inicialmente mobilizada, já que alguns ainda não se conheciam; o
projeto, com as atividades a serem desenvolvidas; o cronograma de
execução das ações de qualificação dos periódicos; e o orçamento
aprovado. Este último sofreu corte na verba de custeio, todavia
preservou todas as bolsas solicitadas, não prejudicando a execução do
projeto e consecutivamente da RIECE.

Contudo, foi na ata da segunda reunião do grupo que foi registrada a
criação oficial da RIECE, em 9 de julho de 2024. Na ocasião, o referido
encontro virtual foi coordenado por Lia Machado Fiuza Fialho, da
Universidade Estadual do Ceará (UECE), e contou com a presença de: Ana
Maria Cetto -- México (Latindex), Nelson Antonio Simâo Gimenes -- PUC/SP
(EDUC@), Sylvia Laura Piovesan Suárez -- Uruguai (AURA), Benjamín Barón
Velandia -- Colômbia, Eva Álvarez Ramos -- Espanha, Christian Andrés
Torres Hurtado -- Colômbia, María Alejandra Balbi -- Uruguai, Eurico
Wongo Gungula -- Angola, Jefferson Mainardes -- UEPG, Charliton José dos
Santos Machado -- UFPB, Fabiane Maia Garcia -- UFAM, Olivia Morais de
Medeiros Neta -- UFRN, Karla Angélica Silva do Nascimento --
UniChristus/UECE, Maria Lília Imbiriba Sousa Colares -- Ufopa, Cláudia
Tavares do Amaral -- UFCAT, Karla Colares Vasconcelos - UFRR, Gilson
Xavier Azevedo -- UEG, Marcelo Silva de Souza Ribeiro -- Univasf, Rafael
Fonseca de Castro -- UNIR, Claudio Nunes -- UESB, Tulio Ramírez --
Venezuela, Jaime Padilla Verdugo -- Equador, Héctor Magaña Vargas --
México, Miguel Ángel Herrera Pavo -- Equador, Abelardo Juvenal Montiel
Benitez -- Paraguai, Pedro José Rivas -- Venezuela e Mariana Lacaio --
Costa Rica. José Airton de Freitas Pontes Junior -- UECE, Cristine
Brandenburg -- UECE, Alfrancio Dias -- UFS, Rebeca Estefano de Salazar
-- Venezuela e Segundo Aguado (Redalic) justificaram ausência.

Como se registra em ata, ``foi aprovada a logomarca da Rede
Ibero-Americana de Editores em Educação (RIECE) e o selo de qualidade
que será concedido às revistas da rede que cumprirem critérios de
qualidade e ciência aberta'' (Ata RIECE, 9 de julho de 2024).

\begin{figure}

\begin{minipage}{0.50\linewidth}

\centering{

\pandocbounded{\includegraphics[keepaspectratio]{img/riece.png}}

}

\subcaption{\label{fig-2}Logomarca da RIECE}

\end{minipage}%
%
\begin{minipage}{0.50\linewidth}

\centering{

\pandocbounded{\includegraphics[keepaspectratio]{img/selo.png}}

}

\subcaption{\label{fig-3}Selo de qualidade}

\end{minipage}%

\caption{\label{fig-1}Logomarca e selo}

\end{figure}%

A logomarca, de fácil visualização e leitura, destaca a letra ``e'',
correspondente à palavra ``editores'', e está acompanhada de sua
interpretação literal escrita em português e inglês, idiomas
predominantemente utilizados nos países ibero-americanos envolvidos no
projeto.

Importa salientar que a Rede valoriza a ciência aberta (Unesco, 2017,
2021) e a rota diamante, inclusive, pelo fato de que não se faz
restrição aos editores de periódicos científicos quanto ao ingresso na
RIECE, todavia, para a concessão do selo, já havia regras nítidas, tais
como um tempo mínimo de 5 anos de existência do periódico, não possuir
características predatórias e adotar a rota diamante. Essa rota é
caracterizada pelo acesso aberto aos textos publicados, distribuídos e
preservados pelos periódicos sem custos para o leitor ou para o autor
(Appel \& Albagli, 2019; Pereira \& Furnival, 2020). Dessa maneira,
apenas as revistas que não cobram o pagamento das taxas de processamento
de artigo podem obter o selo, conhecido como Article Processing Charges
(APCs).

Constatou-se, a partir da criação da RIECE, de maneira concreta, que
houve três ações consideradas como as mais preponderantes em termos de
importância: 1) o diagnóstico dos 32 periódicos envolvidos no projeto,
destacando suas fortalezas e suas fraquezas, registrado no documento
denominado ``Relatório diagnóstico das revistas envolvidas na RIECE'';
2) as missões de internacionalização iniciadas, registradas na
Plataforma Carlos Chagas, site do governo brasileiro onde registra a
execução do projeto; 3) as atividades de participações em eventos e
formações para a qualificação no campo editorial e para a divulgação e
fortalecimento da RIECE, registradas a partir de publicação em anais de
eventos com publicações sobre a RIECE ou com certificados, palestras,
coordenação de mesas, minicursos ou oficinas. Tais ações serão descritas
a seguir de maneira mais detalhada.

\begin{enumerate}
\def\labelenumi{\arabic{enumi}.}
\tightlist
\item
  O relatório diagnóstico dos 32 periódicos envolvidos no projeto,
  destacando suas fortalezas e suas fraquezas
\end{enumerate}

A partir da análise do relatório diagnóstico, foi possível perceber que
15 das revistas envolvidas na RIECE eram relativamente jovens, sendo
criadas entre 2011 e 2020, todavia, existiam outras 15 que já possuíam
mais de 15 anos de existência e apenas três possuíam menos de dois anos
de criação.

\begin{longtable}[]{@{}lllllll@{}}
\caption{Tempo de existência das revistas}\label{tbl-4}\tabularnewline
\toprule\noalign{}
\multirow{2}{=}{Ano} & \multicolumn{2}{l}{%
Brasil} & \multicolumn{2}{l}{%
Demais países} & \multicolumn{2}{l@{}}{%
Total} \\
& n & \% & n & \% & n & \% \\
\midrule\noalign{}
\endfirsthead
\toprule\noalign{}
\multirow{2}{=}{Ano} & \multicolumn{2}{l}{%
Brasil} & \multicolumn{2}{l}{%
Demais países} & \multicolumn{2}{l@{}}{%
Total} \\
& n & \% & n & \% & n & \% \\
\midrule\noalign{}
\endhead
\bottomrule\noalign{}
\endlastfoot
1988-2000 & 3 & 19\% & 4 & 24\% & 7 & 21\% \\
2001-2010 & 4 & 25\% & 4 & 24\% & 8 & 24\% \\
2011-2020 & 7 & 44\% & 8 & 47\% & 15 & 45\% \\
2021-2023 & 2 & 13\% & 1 & 6\% & 3 & 9\% \\
\textbf{Total} & \textbf{16} & \textbf{100\%} & \textbf{17} &
\textbf{100\%} & \textbf{33} & \textbf{100\%} \\
\end{longtable}

No tocante à periodicidade da publicação, chamou a atenção o fato de que
a maioria já adotava o fluxo contínuo, somando-se 20 periódicos. No
entanto, ainda havia 13 revistas que optavam por publicação em números
periódicos -- semestrais, quadrimestrais ou trimestrais.

\begin{longtable}[]{@{}lllllll@{}}
\caption{Periodicidade da publicação}\label{tbl-5}\tabularnewline
\toprule\noalign{}
\multirow{2}{=}{Periodicidade} & \multicolumn{2}{l}{%
Brasil} & \multicolumn{2}{l}{%
Demais países} & \multicolumn{2}{l@{}}{%
Total geral} \\
& n & \% & n & \% & n & \% \\
\midrule\noalign{}
\endfirsthead
\toprule\noalign{}
\multirow{2}{=}{Periodicidade} & \multicolumn{2}{l}{%
Brasil} & \multicolumn{2}{l}{%
Demais países} & \multicolumn{2}{l@{}}{%
Total geral} \\
& n & \% & n & \% & n & \% \\
\midrule\noalign{}
\endhead
\bottomrule\noalign{}
\endlastfoot
Contínuo & 15 & 88\% & 5 & 31\% & 20 & 61\% \\
Trimestral & 0 & 0\% & 1 & 6\% & 1 & 3\% \\
Quadrimestral & 2 & 12\% & 1 & 6\% & 3 & 9\% \\
Semestral & 0 & 0\% & 9 & 56\% & 9 & 27\% \\
\textbf{Total} & \textbf{17} & \textbf{100\%} & \textbf{16} &
\textbf{100\%} & \textbf{33} & \textbf{100\%} \\
\end{longtable}

Também foi dada importância à questão da preservação digital e chamou a
atenção o fato de apenas 17 revistas cuidarem de preservar as
publicações e demais informações relevantes, de modo que 16 revistas
precisam urgentemente cuidar para assegurar a preservação digital. Vale
ressaltar, inclusive, que no Brasil a preservação na Rede Cariniana é
gratuita e os editores pareciam desconhecer essa informação, o que pode
ser facilmente resolvido desde a troca de experiências e saberes em
rede.

\begin{longtable}[]{@{}llll@{}}
\caption{Preservação digital das revistas}\label{tbl-6}\tabularnewline
\toprule\noalign{}
& Brasil & Demais países & Total geral \\
\midrule\noalign{}
\endfirsthead
\toprule\noalign{}
& Brasil & Demais países & Total geral \\
\midrule\noalign{}
\endhead
\bottomrule\noalign{}
\endlastfoot
CLOCKSS & 4 & 1 & 5 \\
LOCKSS & 4 & 4 & 8 \\
Cariniana & 4 & & 4 \\
\textbf{Total} & \textbf{12} & \textbf{5} & \textbf{17} \\
\end{longtable}

No que diz respeito às bases e indexadores, observa-se que há uma
preocupação das revistas em buscar sua divulgação mediante bases e
indexadores diversos.

\begin{longtable}[]{@{}lll@{}}
\caption{Bases e indexadores das revistas}\label{tbl-7}\tabularnewline
\toprule\noalign{}
Indexador & Brasil & Demais países \\
\midrule\noalign{}
\endfirsthead
\toprule\noalign{}
Indexador & Brasil & Demais países \\
\midrule\noalign{}
\endhead
\bottomrule\noalign{}
\endlastfoot
Latindex & 16 & 12 \\
DOAJ & 8 & 11 \\
Dialnet & 5 & 8 \\
Latinrev & 5 & 7 \\
BASE & 7 & 5 \\
ERIH-Plus & 3 & 9 \\
MIAR & 3 & 8 \\
Redalyc & 3 & 7 \\
Scielo & 1 & 8 \\
Road & 7 & 1 \\
Iresie & 3 & 4 \\
Researchbib & 6 & 1 \\
Sherpa Romeo & 3 & 4 \\
CIRC & 1 & 6 \\
Scilit & 3 & 3 \\
Open Aire & 4 & 2 \\
SIS & 3 & 2 \\
MLA & 1 & 4 \\
I2OR & 4 & 1 \\
Latino Americana & 2 & 3 \\
Scopus & 1 & 3 \\
Clase & 3 & 1 \\
Redalyc.org & 1 & 0 \\
IPIndexing & 1 & 0 \\
ESJI & 4 & 0 \\
Amelica & 1 & 3 \\
WoS & 1 & 2 \\
Educ@ & 3 & 0 \\
ERA & 1 & 2 \\
Research4 Life & 3 & 0 \\
Science Gate & 2 & 1 \\
Root Indexing & 2 & 1 \\
FATCAT & 3 & 0 \\
Mir@bel & 2 & 1 \\
EZ3 & 1 & 1 \\
InfoBasa Index & 1 & 1 \\
MAKTABA & 2 & 0 \\
JISC & 2 & 0 \\
ZDB & 2 & 0 \\
Scope Database & 2 & 0 \\
Biblat & 0 & 2 \\
Clacso & 1 & 0 \\
Jisc & 1 & 0 \\
ARDI & 1 & 0 \\
Advenced Sciences Index & 1 & 0 \\
Modern Language Association & 1 & 0 \\
IPN & 1 & 0 \\
Proquest & 1 & 0 \\
\end{longtable}

No quesito uso de programas específicos para a prevenção de
similaridades, em especial, do tipo plágio, surpreendeu o fato de 18 das
revistas não prestarem essa informação aos seus leitores e autores.
Dentre as 15 que disponibilizavam essa informação, o programa mais
utilizado foi o Turnitin.

\begin{longtable}[]{@{}
  >{\raggedright\arraybackslash}p{(\linewidth - 6\tabcolsep) * \real{0.6571}}
  >{\raggedright\arraybackslash}p{(\linewidth - 6\tabcolsep) * \real{0.0762}}
  >{\raggedright\arraybackslash}p{(\linewidth - 6\tabcolsep) * \real{0.1429}}
  >{\raggedright\arraybackslash}p{(\linewidth - 6\tabcolsep) * \real{0.1238}}@{}}
\caption{Utilização de programas para identificação de
similaridade}\label{tbl-8}\tabularnewline
\toprule\noalign{}
\begin{minipage}[b]{\linewidth}\raggedright
\end{minipage} & \begin{minipage}[b]{\linewidth}\raggedright
Brasil
\end{minipage} & \begin{minipage}[b]{\linewidth}\raggedright
Demais países
\end{minipage} & \begin{minipage}[b]{\linewidth}\raggedright
Total geral
\end{minipage} \\
\midrule\noalign{}
\endfirsthead
\toprule\noalign{}
\begin{minipage}[b]{\linewidth}\raggedright
\end{minipage} & \begin{minipage}[b]{\linewidth}\raggedright
Brasil
\end{minipage} & \begin{minipage}[b]{\linewidth}\raggedright
Demais países
\end{minipage} & \begin{minipage}[b]{\linewidth}\raggedright
Total geral
\end{minipage} \\
\midrule\noalign{}
\endhead
\bottomrule\noalign{}
\endlastfoot
Não informado & 10 & 8 & 18 \\
Turnitin & & 6 & 6 \\
Copy Spider & 3 & & 3 \\
Similarity Check & 2 & 1 & 3 \\
Copy Spider, Plagiarism & 1 & & 1 \\
Grammarly, Plagium, Copionic, WriteCheck, PaperRater, Plagarisma.net,
Viper Plagarism Scanner ou CrossCheck & & 1 & 1 \\
iThenticate & 1 & 0 & 1 \\
\textbf{Total} & \textbf{17} & \textbf{16} & \textbf{33} \\
\end{longtable}

Consoante a ética na política editorial das revistas, observou-se que 24
revistas tomavam como base para a elaboração de diretrizes documentos
nacionais ou internacionais amplamente divulgados, majoritariamente
fundamentas no Committee on Publication Ethics (COPE). No entanto, havia
ainda nove revistas que não explicitavam com clareza quais procedimentos
éticos adotavam.

\begin{longtable}[]{@{}
  >{\raggedright\arraybackslash}p{(\linewidth - 6\tabcolsep) * \real{0.6875}}
  >{\raggedright\arraybackslash}p{(\linewidth - 6\tabcolsep) * \real{0.0833}}
  >{\raggedright\arraybackslash}p{(\linewidth - 6\tabcolsep) * \real{0.1562}}
  >{\raggedright\arraybackslash}p{(\linewidth - 6\tabcolsep) * \real{0.0729}}@{}}
\caption{Normas que fundamentam a política
editorial}\label{tbl-9}\tabularnewline
\toprule\noalign{}
\begin{minipage}[b]{\linewidth}\raggedright
Norma
\end{minipage} & \begin{minipage}[b]{\linewidth}\raggedright
Brasil
\end{minipage} & \begin{minipage}[b]{\linewidth}\raggedright
Demais países
\end{minipage} & \begin{minipage}[b]{\linewidth}\raggedright
Total
\end{minipage} \\
\midrule\noalign{}
\endfirsthead
\toprule\noalign{}
\begin{minipage}[b]{\linewidth}\raggedright
Norma
\end{minipage} & \begin{minipage}[b]{\linewidth}\raggedright
Brasil
\end{minipage} & \begin{minipage}[b]{\linewidth}\raggedright
Demais países
\end{minipage} & \begin{minipage}[b]{\linewidth}\raggedright
Total
\end{minipage} \\
\midrule\noalign{}
\endhead
\bottomrule\noalign{}
\endlastfoot
Código de Ética da UTIC & & 1 & 1 \\
COPE (Committee on Publication Ethics) & 9 & 12 & 21 \\
Manual da American Psychological Association - APA & 1 & & 1 \\
SciELO e COPE (Committee on Publication Ethics) & 1 & & 1 \\
Não possui & 6 & 3 & 9 \\
\textbf{Total} & \textbf{17} & \textbf{16} & \textbf{33} \\
\end{longtable}

Outras informações importantes também foram diagnósticas no relatório, a
exemplo da cobrança de APC, exercida por apenas duas revistas, bem como
na internacionalização no corpo editorial e de autores, que demonstrou
que os periódicos que já trabalhavam nessa direção.

Dentre as dificuldades editoriais referentes aos periódicos estritamente
relacionados à internacionalização, podem-se listar, dentre outras: a
dificuldade de compor um comitê científico/editorial internacional com
representação de vários países e continentes; conseguir pareceristas
internacionais capacitados para emitir bons pareceres em tempo hábil;
indexar em bases/diretórios/indexadores internacionais; publicar artigos
produzidos por pesquisadores internacionais; divulgar os artigos
publicados mundo afora e conseguir que a revista seja lida e citada
internacionalmente.

Com o diagnóstico prévio, mostrando as possíveis fragilidades dos
periódicos, o objetivo é investir para sanar essas lacunas com o mote de
qualificar o periódico e o trabalho do editor. Inclusive, a partir de
missões de internacionalização, há a possibilidade de editores de países
distintos interagirem presencialmente, sendo tutoreados por colegas mais
experientes, para investir na melhoria da qualidade dos periódicos da
RIECE.

\begin{enumerate}
\def\labelenumi{\arabic{enumi}.}
\setcounter{enumi}{1}
\tightlist
\item
  As missões de internacionalização
\end{enumerate}

Na ata de criação da RIECE, de 9 de julho de 2024, destacou-se para os
editores internacionais que os brasileiros seriam financiados para
missões de internacionalização e que iriam entrar em contato para
iniciar a organização das missões de internacionalização. Segundo
registro na Plataforma Carlos Chagas, apenas uma missão já havia sido
realizada; a maioria estava agendada para o primeiro semestre de 2025 e
as demais seriam desenvolvidas no segundo semestre.

\begin{longtable}[]{@{}
  >{\raggedright\arraybackslash}p{(\linewidth - 6\tabcolsep) * \real{0.2412}}
  >{\raggedright\arraybackslash}p{(\linewidth - 6\tabcolsep) * \real{0.4353}}
  >{\raggedright\arraybackslash}p{(\linewidth - 6\tabcolsep) * \real{0.2000}}
  >{\raggedright\arraybackslash}p{(\linewidth - 6\tabcolsep) * \real{0.1235}}@{}}
\caption{Editores já cadastrados para missão de
internacionalização}\label{tbl-10}\tabularnewline
\toprule\noalign{}
\begin{minipage}[b]{\linewidth}\raggedright
Bolsista/Editor
\end{minipage} & \begin{minipage}[b]{\linewidth}\raggedright
Processo
\end{minipage} & \begin{minipage}[b]{\linewidth}\raggedright
Vigência
\end{minipage} & \begin{minipage}[b]{\linewidth}\raggedright
Categoria da bolsa
\end{minipage} \\
\midrule\noalign{}
\endfirsthead
\toprule\noalign{}
\begin{minipage}[b]{\linewidth}\raggedright
Bolsista/Editor
\end{minipage} & \begin{minipage}[b]{\linewidth}\raggedright
Processo
\end{minipage} & \begin{minipage}[b]{\linewidth}\raggedright
Vigência
\end{minipage} & \begin{minipage}[b]{\linewidth}\raggedright
Categoria da bolsa
\end{minipage} \\
\midrule\noalign{}
\endhead
\bottomrule\noalign{}
\endlastfoot
Charliton José dos Santos Machado &
\href{http://efomento.cnpq.br/efomento/contrato/indicacaoBolsista.do?acao=detalharIndicacao&codigoSolicitacao=4810698&codigoProcesso=20234424339}{201545/2024-1}
& 01/01/2025 a 31/01/2025 & DES \\
Fabiane Maia Garcia &
\href{http://efomento.cnpq.br/efomento/contrato/indicacaoBolsista.do?acao=detalharIndicacao&codigoSolicitacao=4810698&codigoProcesso=20234424339}{201579/2024-3}
& 01/02/2025 a 28/02/2025 & DES \\
Lia Machado Fiuza Fialho &
\href{http://efomento.cnpq.br/efomento/contrato/indicacaoBolsista.do?acao=detalharIndicacao&codigoSolicitacao=4810698&codigoProcesso=20234424339}{201538/2024-5}
& 01/03/2025 a 31/03/2025 & DES \\
Maria Aparecida Alves da Costa &
\href{http://efomento.cnpq.br/efomento/contrato/indicacaoBolsista.do?acao=detalharIndicacao&codigoSolicitacao=4810698&codigoProcesso=20234424339}{201565/2024-2}
& 01/01/2025 a 31/01/2025 & DEJ \\
Maria Lília Imbiriba Sousa Colares &
\href{http://efomento.cnpq.br/efomento/contrato/indicacaoBolsista.do?acao=detalharIndicacao&codigoSolicitacao=4810698&codigoProcesso=20234424339}{201646/2024-2}
& 01/02/2025 a 28/02/2025 & DES \\
Olivia Morais de Medeiros Neta &
\href{http://efomento.cnpq.br/efomento/contrato/indicacaoBolsista.do?acao=detalharIndicacao&codigoSolicitacao=4810698&codigoProcesso=20234424339}{201578/2024-7}
& 01/02/2025 a 28/02/2025 & DES \\
Rafael Fonseca de Castro &
\href{http://efomento.cnpq.br/efomento/contrato/indicacaoBolsista.do?acao=detalharIndicacao&codigoSolicitacao=4810698&codigoProcesso=20234424339}{201439/2024-7}
& 01/10/2024 a 31/10/2024 & DES \\
\end{longtable}

Apenas Rafael Fonseca de Castro já havia concluído a missão; outros seis
editores já estavam com suas missões agendadas e os demais estavam em
fase de negociação com as suas instituições de vínculo empregatício e
com o colega que iria trabalhar conjuntamente para qualificar as suas
revistas. Dessa maneira, ainda que o registro da atividade seja muito
importante por ser uma iniciativa inovadora, ainda não havia resultados
que pudessem ser devidamente analisados.

\begin{enumerate}
\def\labelenumi{\arabic{enumi}.}
\setcounter{enumi}{2}
\tightlist
\item
  As atividades de participações em eventos e formações para a
  qualificação no campo editorial e para a divulgação e fortalecimento
  da RIECE
\end{enumerate}

Mesmo com o tempo diminuto de criação da RIECE, menos de um ano, é
possível observar uma profícua mobilização dos envolvidos no projeto
para trabalhar em prol da articulação e divulgação da Rede. Dentre as
atividades mais expressivas nessa direção, mencionam-se:

\begin{quote}
Participação na 15ª Conferência Lusófona de Ciência Aberta (ConfOA), com
a temática ``Acesso aberto e dados de investigação abertos: sistemas,
políticas e práticas'', em que apresentaram trabalho denominado ``Rede
Ibero-Americana de Editores Científicos de Educação (RIECE):
internacionalização, indexação e aderência à ciência aberta'' e o
trabalho ``O descompasso da emergência da Ciência Aberta com os
critérios Qualis Capes Educação''. Na ocasião, questionava-se a forma de
avaliar a produção e os periódicos baseando-se majoritariamente em
métricas e se divulgava a RIECE convidando editores da área de Educação
a ingressarem na Rede e robustecer o debate (Fialho et~al., 2024).
\end{quote}

\begin{quote}
Sessão de comunicação oral realizada no congresso ABEC Meeting 2023,
sobre ``A importância e os desafios da avaliação aberta à luz da
experiência da Revista Práticas Educativas, Memórias e Oralidades'', na
qual os editores eram incentivados a adotarem a avaliação aberta
consoante os princípios de transparência da ciência aberta (Nascimento
\& Fialho, 2023).
\end{quote}

\begin{quote}
Palestra no Encontro do Fórum de Editores, com o título ``Desafios e
possibilidades na editoria de revistas científicas'', durante o XII
Congresso ABRACE, realizado em junho de 2023 na Universidade Federal do
Pará, com o mote de debater boas práticas editoriais, o alinhamento à
ciência aberta e divulgar a RIECE. Apresentação de trabalho no ``Eixo
Temático Avaliação de Periódicos da Área de Educação'', no IV Congresso
Nacional de Editores de Periódicos de Educação (Coneped) -- Fepae/Anped,
evento presencial realizado em maio de 2024 na Universidade do Estado do
Rio de Janeiro, ocasião em que se defendia uma avaliação que
considerasse as especificidades do sul global e fortalecesse a ciência
em países latino-americanos, africanos e caribenhos, bem como
fortalecesse a RIECE. Participação da coordenadora na mesa-redonda ``A
consolidação e os desafios de periódicos científicos cearenses'', na
condição de palestrante, promovida pela Revista Ensino em Debate em maio
de 2024, com o objetivo de envolver os editores da Rede no debate acerca
das problemáticas que precisam ser enfrentadas com vistas à qualificação
e consolidação das revistas científicas de educação da Iberoamérica.
Desenvolvimento de oficina para a indexação de revistas científicas,
realizadas on-line, com o mote de preparar as revistas para atender aos
critérios dos indexadores nacionais e internacionais. Minicurso sobre as
ferramentas disponíveis no OJS, versão 3, e sua utilização qualificada
para customizar, instalar plugins, preencher metadados e gerar
relatórios com o mote de otimizar a exploração de recursos disponíveis
na plataforma (Shintaku et~al., 2014).
\end{quote}

Cabe destacar que em todas as atividades era disponibilizado o QR Code
que direcionava ao questionário on-line para que outros editores de
Educação pudessem se cadastrar na RIECE, qual seja:

\begin{figure}

\centering{

\pandocbounded{\includegraphics[keepaspectratio]{img/qr.png}}

}

\caption{\label{fig-4}QR Code para juntar-se à RIECE}

\end{figure}%

Desse modo, observa-se que o objetivo é ampliar a Rede inicialmente
constituída e poder constituir uma ampla articulação internacional de
editores da área de Educação que queiram investir esforços em qualificar
sua atuação profissional e os periódicos em que atuam como editores.

\section{Considerações finais}\label{considerauxe7uxf5es-finais}

A relevância do projeto ganhou força ao se constar que são escassas as
revistas educacionais latino-americanas que conseguem compor o catálogo
dos grandes sponsors e internacionalizarem-se como periódicos de
referência para o mundo. Diante disso, elaborou-se o pressuposto de que
a constituição de uma rede de cooperação poderia fomentar a
internacionalização e o desenvolvimento científico e tecnológico no
âmbito da divulgação científica, desde a socialização de conhecimentos e
ajuda mútua, tutoreados por editores de revistas mais consolidadas.

Essa articulação internacional permitiria discutir e implementar
inovações no campo editorial, de modo que o avanço tecnológico
possibilitará, além de ampliar a internacionalização, investir na
qualificação das revistas científicas envolvidas no projeto, o que
reflete no aumento da visibilidade e da credibilidade dos periódicos e
na disseminação das pesquisas científicas com maior alinhamento à
ciência aberta e às melhores práticas editoriais mundiais (Fialho,
2023). Afinal, a popularização do conhecimento, socializado amplamente
de maneira democrática, em acesso aberto, e a disseminação qualificada
do conhecimento repercutem no desenvolvimento científico e tecnológico
do país e na almejada internacionalização (Kunsch, 2004).

Conclui-se que a RIECE vem colaborando para qualificar científica e
tecnologicamente os periódicos de Educação a partir da articulação dos
conhecimentos editoriais da equipe compartilhados voluntariamente,
consoante a ciência aberta, entre os seus integrantes. Sua principal
meta é ampliar o número de países da Rede e aumentar os editores de
Educação comprometidos com a proposta.

A cooperação internacional proposta mediante a criação da RIECE tem se
mostrado de suma relevância, porque permite, a partir da articulação
entre os editores, a ampliação de membros internacionais atuando nos
conselhos editoriais e científicos dos periódicos brasileiros, e
vice-versa, a troca de indicação de pareceristas internacionais, a
colaboração para a divulgação da revista em outros países, a captação de
textos internacionais e a tutoria com editores mais experientes para
buscar adequação à ciência aberta, aos critérios dos indexados para
ampliar a indexação, dentre outros intercâmbios desde o estreitamento de
relações com pesquisadores estrangeiros.

\bookmarksetup{startatroot}

\chapter{Summary}\label{summary}

In summary, this book has no content whatsoever.

\chapter*{Referencias}\label{bibliography}
\addcontentsline{toc}{chapter}{Referencias}

\phantomsection\label{refs}
\begin{CSLReferences}{1}{0}
\bibitem[\citeproctext]{ref-Anna2019}
Anna, J. S. (2019). Comunicação científica e o papel dos periódicos
científicos no desenvolvimento das ciências. \emph{Biblionline},
\emph{15}(1), 3-18.
\url{https://doi.org/10.22478/ufpb.1809-4775.2019v15n1.44365}

\bibitem[\citeproctext]{ref-Appel2019}
Appel, A. L., \& Albagli, S. (2019). The adoption of Article Processing
Charges as a business model by Brazilian Open Access journals.
\emph{Transinformação}, \emph{31}, e180045.
\url{https://doi.org/10.1590/2318-0889201931e180045}

\bibitem[\citeproctext]{ref-Camargo2018}
Camargo, B. V., \& Justo, A. M. (2018). \emph{Tutorial para uso do
software {[}de análise textual{]} Iramuteq}. UFSC.

\bibitem[\citeproctext]{ref-Fialho2023}
Fialho, L. M. F. (2023). \emph{Ciência Aberta, uma compreensão inicial}.
Comissão de ética e pesquisa da ANPEd.

\bibitem[\citeproctext]{ref-Fialho2024}
Fialho, L. M. F., Nascimento, K. A. S., \& Costa, M. A. A. (2024). Rede
Ibero-Americana de Editores Científicos de Educação (RIECE):
internacionalização, indexação e aderência à ciência aberta. \emph{15ª
Conferência Lusófona de Ciência Aberta (ConfOA)}, \emph{53}.

\bibitem[\citeproctext]{ref-Fontes2021}
Fontes, I. (2021). \emph{Competências do editor-chefe de periódico
científico: gaps e tendências} {[}Mathesis{]}. Fundação Getúlio Vargas.

\bibitem[\citeproctext]{ref-Gomes2010}
Gomes, V. P. (2010). O editor de revista científica: desafios da prática
e da formação. \emph{Informação \& Informação}, \emph{15}(1), 147-172.

\bibitem[\citeproctext]{ref-Kunsch2004}
Kunsch, W. L. (2004). O editor científico. \emph{Revista Acadêmica do
Grupo Comunicacional de São Bernardo}, \emph{1}(1).

\bibitem[\citeproctext]{ref-Ludke1986}
Lüdke, M., \& André, M. E. D. A. (1986). \emph{Pesquisa em educação:
abordagens qualitativas}. EPU.

\bibitem[\citeproctext]{ref-Minayo2008}
Minayo, M. C. S. (2008). \emph{O desafio do conhecimento: pesquisa
qualitativa em saúde}. Hucitec.

\bibitem[\citeproctext]{ref-Morais2022}
Morais, A., Rode, S. M., \& Galleti, S. (2022). \emph{Desafios e
perspectivas da editoria científica: memórias críticas do ABEC Meeting
Live 2021}. ABEC Brasil.

\bibitem[\citeproctext]{ref-Nascimento2023}
Nascimento, K. A. S., \& Fialho, L. M. F. (2023). A importância e os
desafios da avaliação aberta à luz da experiência da Revista Práticas
Educativas, Memórias e Oralidades. \emph{Abec Meeting}.

\bibitem[\citeproctext]{ref-Oliveira2020}
Oliveira, M. M. (2020). \emph{Como fazer pesquisa qualitativa}. Vozes.

\bibitem[\citeproctext]{ref-Pereira2020}
Pereira, V., \& Furnival, A. C. (2020). Revistas científicas em acesso
aberto brasileiras no DOAJ: modelos de negócio e sua sustentabilidade
financeira. \emph{Brazilian Journal of Information Science: Research
Trends}, \emph{14}(1), 88-111.
\url{https://doi.org/10.36311/1981-1640.2020.v14n1.05.p88}

\bibitem[\citeproctext]{ref-Pimentel2001}
Pimentel, A. (2001). O método da análise documental: seu uso numa
pesquisa histórica. \emph{Cadernos de Pesquisa}, \emph{114}, 179-195.

\bibitem[\citeproctext]{ref-SaSilva2009}
Sá-Silva, J. R., Almeida, C. D., \& Guindani, J. F. (2009). Pesquisa
documental: pistas teóricas e metodológicas. \emph{Revista Brasileira de
História \& Ciências Sociais}, \emph{1}(1).

\bibitem[\citeproctext]{ref-Shintaku2014}
Shintaku, M. et~al. (2014). \emph{Cartilha para adequação de publicações
SEER/OJS ao Latindex}. Instituto Brasileiro de Informação em Ciência e
Tecnologia.

\bibitem[\citeproctext]{ref-Tripp2005}
Tripp, D. (2005). Pesquisa-ação: uma introdução metodológica.
\emph{Educação e Pesquisa}, \emph{31}(3), 443-466.
\url{https://doi.org/10.1590/S1517-97022005000300009}

\bibitem[\citeproctext]{ref-Unesco2017}
Unesco. (2017). \emph{Recomendación sobre la Ciencia y los
Investigadores Científicos}. Unesco.

\bibitem[\citeproctext]{ref-Unesco2021}
Unesco. (2021). \emph{Recomendação da Unesco sobre Ciência Aberta}.
Unesco.

\bibitem[\citeproctext]{ref-Werlang2021}
Werlang, E., Carlim, P. E., \& Rode, S. M. (2021). \emph{Desafios e
perspectivas da editoria científica: memórias críticas do ABEC Meeting
Live 2020}. ABEC Brasil.

\bibitem[\citeproctext]{ref-Yin2005}
Yin, R. K. (2005). \emph{Estudo de caso: planejamento e métodos}.
Bookman.

\end{CSLReferences}


\backmatter


\end{document}
